\thispagestyle{empty}
\frontmatter
\begin{titlepage}
\vspace{5mm}
\begin{figure}[hbtp]
\centering
\includegraphics[scale=.13]{Images/UNIPD.png}
\end{figure}
\vspace{5mm}
\begin{center}
{{\huge{\textsc{\bf UNIVERSIT\`A DEGLI STUDI DI PADOVA}}}\\}
\vspace{5mm}
\hrule
\vspace{5mm}
{\Large{\bf Dipartimento di Fisica e Astronomia ``Galileo Galilei''}} \\
\vspace{5mm}
{\Large{\textsc{\bf Corso di Laurea in Fisica}}}\\
\vspace{20mm}
{\Large{\textsc{\bf Tesi di Laurea}}}\\
\vspace{20mm}
\begin{spacing}{2}
{\LARGE \textbf{Development of a ML-based model-independent \\analysis strategy at the LHC}}\\
\end{spacing}
\vspace{8mm}
\end{center}

\vspace{10mm}
%----------------------------------------------------------------------

\begin{spacing}{2}
\begin{tabular}{ l  c  c c c  cccccccc c c c  c  l }
{\Large{\bf Relatore:}} &&&&&&&&&&&&&&&&& {\Large{\bf Laureando:}}\\
{\Large{\bf \textit{Prof. Marco Zanetti}}} &&&&&&&&&&&&&&&&& {\Large{\bf \textit{Rocco Ardino}}}\\
% {\Large{\bf Correlatore}}\\
% {\Large{\bf Prof./Dr. Nome Cognome}}\\
\end{tabular}
\end{spacing}

%----------------------------------------------------------------------
\begin{comment}
\begin{spacing}{2}
\begin{flushleft}                           
{\Large \textbf{Relatore\colon}	\\
\textbf{\textit{Prof. Marco Zanetti}}} \\
{\Large \textbf{Laureando\colon}	\\
\textbf{\textit{Rocco Ardino}}}
\end{flushleft}
\end{spacing}
\end{comment}
\vspace{15 mm}
\vspace*{\fill}
\hrule
\begin{center}
{\Large{\bf Anno Accademico 2018/2019}}
\end{center}
\end{titlepage}



\vspace*{\fill}
\begingroup

\begin{center}
{{\large{\bf Abstract}}}
\end{center}
%\vspace{10mm}

\begin{center}
Searching for New Physics is the primary goal of the CMS experiment at the LHC. Performing such search without relying on a
specific theory extending the Standard Model (SM) is of paramount importance but at the same time highly non-trivial. Recently,
several proposal have been made in that sense, in particular exploiting the power of modern ML techniques. The goal of this thesis is
to apply such model-independent strategies to concrete physics cases, aiming at detecting signals solely by comparing the collision
data with what predicted by the SM.
\end{center}


\vspace{40mm}


\begin{center}
{{\large{\bf Sommario}}}
\end{center}
%\vspace{10mm}

\begin{center}
La ricerca di nuova fisica è l'obiettivo primario dell'esperimento CMS a LHC. Eseguire una ricerca del genere senza fare affidamento su una teoria specifica che allarga il Modello Standard (SM) è di primaria importanza, ma allo stesso tempo anche altamente non banale. Recentemente, sono state avanzate diverse proposte verso questa direzione, in particolare sfruttando la potenza delle tecniche moderne di Machine Learning. L'obiettivo di questa tesi è applicare tali strategie "model-independent" a casi fisici concreti, mirando a rivelare segnali esclusivamente comparando i dati delle collisioni con quelli predetti dallo SM.
\end{center}

\endgroup
\vspace*{\fill}




\clearpage{\pagestyle{empty}\cleardoublepage}